% ---------------------------------------------------------
% TABELA EMENTA 10: Espanhol — Ano 1
% ---------------------------------------------------------
\begin{xltabular}{\linewidth}{|X|p{2.5cm}|p{2.5cm}|}
\hline
\multirow{3}{=}{\textcolor{ifscgreen}{\textbf{Unidade Curricular:}}\\[3pt]\textbf{\large Espanhol — Ano 1}} & \multicolumn{2}{p{\dimexpr5cm+2\tabcolsep+\arrayrulewidth\relax}|}{\textcolor{ifscgreen}{\textbf{Semestre:}} \textbf{1º e 2º}} \\ \cline{2-3}
 & \textcolor{ifscgreen}{\textbf{CH EaD*:}} & \textcolor{ifscgreen}{\textbf{CH Total*:}} \\ \cline{2-3}
 & \textbf{00 h} & \textbf{80 h} \\ \hline
\endhead
\multicolumn{3}{|p{\dimexpr\linewidth-2\tabcolsep-2\arrayrulewidth\relax}|}{\textcolor{ifscgreen}{\textbf{Objetivos:}}} \\ \hline
\multicolumn{3}{|p{\dimexpr\linewidth-2\tabcolsep-2\arrayrulewidth\relax}|}{
\begin{itemize}[leftmargin=*,noitemsep,topsep=2pt]
 \item Compreender os aspectos fonéticos e de pronúncia introdutórios da língua espanhola.
 \item Aplicar cumprimentos, apresentações e vocabulário básico (cores, números, família, dias da semana) em interações cotidianas.
 \item Utilizar estruturas gramaticais básicas no tempo presente, pronomes pessoais e artigos (definidos e indefinidos) em contextos de comunicação simples.
 \item Comunicar-se em situações práticas do dia a dia, tais como compras, restaurantes e transportes.
 \item Reconhecer e discutir aspectos de interculturalidade, incluindo costumes, tradições, cidadania, segurança no trânsito, práticas sustentáveis (meio ambiente) e direitos humanos em culturas hispanofalantes.
\end{itemize}
} \\ \hline
\multicolumn{3}{|p{\dimexpr\linewidth-2\tabcolsep-2\arrayrulewidth\relax}|}{\textcolor{ifscgreen}{\textbf{Conteúdos:}}} \\ \hline
\multicolumn{3}{|p{\dimexpr\linewidth-2\tabcolsep-2\arrayrulewidth\relax}|}{
\begin{itemize}[leftmargin=*,noitemsep,topsep=2pt]
 \item Introdução à língua espanhola: fonética e pronúncia. Cumprimentos e apresentações. Vocabulário básico: cores, números, família, dias da semana. Estruturas gramaticais básicas: verbos no presente, pronomes pessoais, artigos definidos e indefinidos. Interculturalidade: costumes e tradições em países de língua espanhola. Comunicação em situações cotidianas: compras, restaurantes, transportes. Temáticas transversais: cidadania e segurança no trânsito, práticas sustentáveis em relação ao meio ambiente, e respeito aos direitos humanos em diferentes culturas hispanofalantes.
\end{itemize}
} \\ \hline
\multicolumn{3}{|p{\dimexpr\linewidth-2\tabcolsep-2\arrayrulewidth\relax}|}{\textcolor{ifscgreen}{\textbf{Estratégias de Ensino e Aprendizagem:}}} \\ \hline
\multicolumn{3}{|p{\dimexpr\linewidth-2\tabcolsep-2\arrayrulewidth\relax}|}{- **Metodologia:** Aulas expositivas com uso de recursos audiovisuais (vídeos, áudios) e atividades práticas (diálogos, role-playing); estratégias de ensino colaborativo, promovendo discussões interculturais e resolução de problemas. Em caso de carga horária EAD: utilização exclusiva do Ambiente Virtual de Aprendizagem Moodle, com livros digitais, fóruns de discussão, quizzes e vídeo-aulas; interações síncronas via videoconferências e fóruns de debate, com encontros presenciais agendados para avaliações orais e escritas; tutoria realizada pelo próprio docente da unidade curricular. Sem divisão de turma.
- **Avaliação:** Avaliação contínua baseada em tarefas orais e escritas, além de testes de compreensão auditiva.} \\ \hline
\multicolumn{3}{|p{\dimexpr\linewidth-2\tabcolsep-2\arrayrulewidth\relax}|}{\textcolor{ifscgreen}{\textbf{Bibliografia Básica:}}} \\ \hline
\multicolumn{3}{|p{\dimexpr\linewidth-2\tabcolsep-2\arrayrulewidth\relax}|}{1. FANJUL, Adrián (org.). \textbf{Gramática y práctica de español para brasileños}. 2. ed. São Paulo: Moderna, 2011.
2. MORENO, Concha; TUTS, Martina. \textbf{Cinco estrellas: español para el turismo}. 2. ed. Madrid: SGEL, 2011.} \\ \hline
\multicolumn{3}{|p{\dimexpr\linewidth-2\tabcolsep-2\arrayrulewidth\relax}|}{\textcolor{ifscgreen}{\textbf{Bibliografia Complementar:}}} \\ \hline
\multicolumn{3}{|p{\dimexpr\linewidth-2\tabcolsep-2\arrayrulewidth\relax}|}{1. GONZÁLEZ HERMOSO, Alfredo. \textbf{Conjugar: verbos de España y de América}. Madrid: Edelsa Grupo Didascalia, 2011.
2. HERMOSO, A. G. \textbf{Conjugar es fácil}. [s.l.]: EDELSA, 1997.
3. PEREZ, Aquilino Sanches. \textbf{Diccionario básico de la lengua española}. [s.l.]: SGEL, 1987.
4. BALLESTERO-ALVAREZ, M. E.; BALBÁS, M. S. \textbf{Minidicionário: espanhol-português, português-espanhol}. São Paulo: FTD, 2007.} \\ \hline
\end{xltabular}

\vspace{0.4cm}


% ---------------------------------------------------------
% TABELA EMENTA 11: Espanhol — Ano 2
% ---------------------------------------------------------
\begin{xltabular}{\linewidth}{|X|p{2.5cm}|p{2.5cm}|}
\hline
\multirow{3}{=}{\textcolor{ifscgreen}{\textbf{Unidade Curricular:}}\\[3pt]\textbf{\large Espanhol — Ano 2}} & \multicolumn{2}{p{\dimexpr5cm+2\tabcolsep+\arrayrulewidth\relax}|}{\textcolor{ifscgreen}{\textbf{Semestre:}} \textbf{3º e 4º}} \\ \cline{2-3}
 & \textcolor{ifscgreen}{\textbf{CH EaD*:}} & \textcolor{ifscgreen}{\textbf{CH Total*:}} \\ \cline{2-3}
 & \textbf{00 h} & \textbf{80 h} \\ \hline
\endhead
\multicolumn{3}{|p{\dimexpr\linewidth-2\tabcolsep-2\arrayrulewidth\relax}|}{\textcolor{ifscgreen}{\textbf{Objetivos:}}} \\ \hline
\multicolumn{3}{|p{\dimexpr\linewidth-2\tabcolsep-2\arrayrulewidth\relax}|}{
\begin{itemize}[leftmargin=*,noitemsep,topsep=2pt]
 \item Revisar e consolidar as estruturas gramaticais básicas da língua espanhola.
 \item Compreender e aplicar tempos verbais no passado (pretérito perfeito e imperfeito).
 \item Expandir o vocabulário em eixos temáticos específicos, tais como alimentos, roupas, lazer e clima.
 \item Desenvolver habilidades de comunicação em situações de nível intermediário, capacitando para a descrição de pessoas, lugares e eventos.
 \item Redigir textos breves e funcionais, como cartas e e-mails informais.
 \item Analisar criticamente manifestações da cultura popular e celebrações hispânicas, articulando reflexões sobre direitos humanos, alimentação saudável e História e Cultura Afro-Brasileira e Indígena.
\end{itemize}
} \\ \hline
\multicolumn{3}{|p{\dimexpr\linewidth-2\tabcolsep-2\arrayrulewidth\relax}|}{\textcolor{ifscgreen}{\textbf{Conteúdos:}}} \\ \hline
\multicolumn{3}{|p{\dimexpr\linewidth-2\tabcolsep-2\arrayrulewidth\relax}|}{
\begin{itemize}[leftmargin=*,noitemsep,topsep=2pt]
 \item Revisão de estruturas gramaticais básicas. Tempos verbais no passado: pretérito perfeito e imperfeito. Vocabulário ampliado: alimentos, roupas, lazer, clima. Interculturalidade: celebrações, festividades e cultura popular nos países hispanofalantes. Comunicação em situações intermediárias: descrevendo pessoas, lugares e eventos. Introdução à escrita de textos breves: cartas e e-mails informais. Temáticas transversais: História e Cultura Afro-Brasileira e Indígena em contextos hispanofalantes, alimentação saudável e educação alimentar, e reflexões integradas sobre os direitos humanos.
\end{itemize}
} \\ \hline
\multicolumn{3}{|p{\dimexpr\linewidth-2\tabcolsep-2\arrayrulewidth\relax}|}{\textcolor{ifscgreen}{\textbf{Estratégias de Ensino e Aprendizagem:}}} \\ \hline
\multicolumn{3}{|p{\dimexpr\linewidth-2\tabcolsep-2\arrayrulewidth\relax}|}{- **Metodologia:** Abordagem comunicativa com ênfase em situações da vida real, priorizando o uso prático do idioma por meio de simulações e diálogos; uso de plataformas digitais para o desenvolvimento de atividades de escrita colaborativa e debates sobre temas culturais e transversais. Em caso de oferta de carga horária EAD: utilização exclusiva do Ambiente Virtual de Aprendizagem Moodle, incorporando quizzes interativos, fóruns de discussão orientados e vídeo-aulas; sessões síncronas via videoconferência voltadas à prática de fala e produção escrita, com momentos assíncronos e encontros presenciais opcionais para reforço e avaliações; tutoria realizada diretamente pelo docente da UC. Sem divisão de turma.
- **Avaliação:** Não especificado no documento fonte.} \\ \hline
\multicolumn{3}{|p{\dimexpr\linewidth-2\tabcolsep-2\arrayrulewidth\relax}|}{\textcolor{ifscgreen}{\textbf{Bibliografia Básica:}}} \\ \hline
\multicolumn{3}{|p{\dimexpr\linewidth-2\tabcolsep-2\arrayrulewidth\relax}|}{1. FANJUL, Adrián (org.). \textbf{Gramática y práctica de español para brasileños}. 2. ed. São Paulo: Moderna, 2011.
2. MORENO, Concha; TUTS, Martina. \textbf{Cinco estrellas: español para el turismo}. 2. ed. Madrid: SGEL, 2011.} \\ \hline
\multicolumn{3}{|p{\dimexpr\linewidth-2\tabcolsep-2\arrayrulewidth\relax}|}{\textcolor{ifscgreen}{\textbf{Bibliografia Complementar:}}} \\ \hline
\multicolumn{3}{|p{\dimexpr\linewidth-2\tabcolsep-2\arrayrulewidth\relax}|}{1. GONZÁLEZ HERMOSO, Alfredo. \textbf{Conjugar: verbos de España y de América}. Madrid: Edelsa Grupo Didascalia, 2011.
2. COTO BAUTISTA, Vanessa; TURZA FERRÉ, Anna. \textbf{Tema a tema B1: español lengua extranjera: curso de conversación}. Madrid: Edelsa, 2011.
3. WILDNER, Ana Kaciara; OLIVEIRA, Leandra Cristina de; SOBOTTKA, Mary Anne Warken. \textbf{Espanhol para o turismo}. Florianópolis: Publicação do IFSC, 2014.} \\ \hline
\end{xltabular}

\vspace{0.4cm}


% ---------------------------------------------------------
% TABELA EMENTA 12: Biologia — Ano 1
% ---------------------------------------------------------
\begin{xltabular}{\linewidth}{|X|p{2.5cm}|p{2.5cm}|}
\hline
\multirow{3}{=}{\textcolor{ifscgreen}{\textbf{Unidade Curricular:}}\\[3pt]\textbf{\large Biologia — Ano 1}} & \multicolumn{2}{p{\dimexpr5cm+2\tabcolsep+\arrayrulewidth\relax}|}{\textcolor{ifscgreen}{\textbf{Semestre:}} \textbf{1º e 2º}} \\ \cline{2-3}
 & \textcolor{ifscgreen}{\textbf{CH EaD*:}} & \textcolor{ifscgreen}{\textbf{CH Total*:}} \\ \cline{2-3}
 & \textbf{00 h} & \textbf{80 h} \\ \hline
\endhead
\multicolumn{3}{|p{\dimexpr\linewidth-2\tabcolsep-2\arrayrulewidth\relax}|}{\textcolor{ifscgreen}{\textbf{Objetivos:}}} \\ \hline
\multicolumn{3}{|p{\dimexpr\linewidth-2\tabcolsep-2\arrayrulewidth\relax}|}{
\begin{itemize}[leftmargin=*,noitemsep,topsep=2pt]
 \item Compreender a Biologia como uma ciência em permanente construção, reconhecendo sua importância para o entendimento da origem, evolução e diversificação da vida, bem como sua relação com o desenvolvimento científico, tecnológico e social.
 \item Reconhecer a célula como unidade estrutural, funcional e genética dos seres vivos, compreendendo a organização celular, os processos metabólicos e os mecanismos envolvidos na manutenção da vida.
 \item Entender os processos de multiplicação celular, reprodução sexuada e anatomofisiologia do corpo humano, relacionando-os à continuidade da vida, à diversidade genética e à saúde humana.
 \item Desenvolver habilidades de observação, investigação e experimentação, utilizando procedimentos próprios das ciências biológicas para formular hipóteses, analisar evidências, interpretar resultados e resolver problemas.
 \item Promover o autocuidado, o empoderamento e o respeito à pluralidade de corpos e gêneros por meio da compreensão dos princípios da educação alimentar e nutricional, da saúde integral e da sexualidade na adolescência.
\end{itemize}
} \\ \hline
\multicolumn{3}{|p{\dimexpr\linewidth-2\tabcolsep-2\arrayrulewidth\relax}|}{\textcolor{ifscgreen}{\textbf{Conteúdos:}}} \\ \hline
\multicolumn{3}{|p{\dimexpr\linewidth-2\tabcolsep-2\arrayrulewidth\relax}|}{
\begin{itemize}[leftmargin=*,noitemsep,topsep=2pt]
 \item Origem da vida: teorias e descobertas recentes; astrobiologia.
 \item A invenção do microscópio e a descoberta da célula.
 \item Células procarióticas e eucarióticas (animais e vegetais).
 \item Componentes químicos celulares: água, macromoléculas, vitaminas e sais minerais; Educação Alimentar e Nutricional.
 \item Estrutura e funcionamento celular: membrana plasmática, citoplasma, organelas citoplasmáticas, citoesqueleto e núcleo.
 \item Metabolismo energético: respiração celular, fermentação, fotossíntese e quimiossíntese.
 \item Ciclo celular: interfase e divisão celular (mitose e meiose).
 \item Reprodução animal: fecundação e desenvolvimento embrionário.
 \item Histologia, anatomia e fisiologia humana: tecidos básicos e sistemas orgânicos.
 \item Saúde e sexualidade na adolescência.
\end{itemize}
} \\ \hline
\multicolumn{3}{|p{\dimexpr\linewidth-2\tabcolsep-2\arrayrulewidth\relax}|}{\textcolor{ifscgreen}{\textbf{Estratégias de Ensino e Aprendizagem:}}} \\ \hline
\multicolumn{3}{|p{\dimexpr\linewidth-2\tabcolsep-2\arrayrulewidth\relax}|}{- **Metodologia:** Abordagem dialógica, contextualizada e investigativa, com levantamento de saberes prévios e problematização de situações reais. Aulas expositivo-dialogadas com quadro e projetor multimídia; atividades práticas em campo e laboratório; análise de experimentos e estudos de caso; discussão de textos, imagens e vídeos científicos; uso de atlas digitais, softwares de animação e simulação, modelos didáticos tridimensionais e jogos educativos; debate de temas transversais; visitas técnicas e eventos científicos. Aulas práticas realizadas preferencialmente no pátio do câmpus e no Laboratório de Biociências (média de 20h anuais). SIGAA utilizado como Ambiente Virtual de Aprendizagem.
- **Avaliação:** Perspectiva processual e formativa, observando a evolução dos conhecimentos ao longo do período letivo, por meio de estudos dirigidos; listas de exercícios focados em ENEM e vestibulares; provas teóricas; relatórios de atividades experimentais; projetos, seminários ou oficinas temáticas; e avaliação atitudinal.} \\ \hline
\multicolumn{3}{|p{\dimexpr\linewidth-2\tabcolsep-2\arrayrulewidth\relax}|}{\textcolor{ifscgreen}{\textbf{Bibliografia Básica:}}} \\ \hline
\multicolumn{3}{|p{\dimexpr\linewidth-2\tabcolsep-2\arrayrulewidth\relax}|}{1. JUNQUEIRA, Luiz Carlos Uchoa; CARNEIRO, José. \textbf{Biologia celular e molecular}. 9. ed. Rio de Janeiro: Guanabara Koogan, 2012.
2. SADAVA, David et al. \textbf{Vida: a ciência da biologia: volume 1: célula e hereditariedade}. 8. ed. Porto Alegre: Artmed, 2009.
3. Livro didático fornecido pelo Fundo Nacional de Desenvolvimento da Educação (FNDE).} \\ \hline
\multicolumn{3}{|p{\dimexpr\linewidth-2\tabcolsep-2\arrayrulewidth\relax}|}{\textcolor{ifscgreen}{\textbf{Bibliografia Complementar:}}} \\ \hline
\multicolumn{3}{|p{\dimexpr\linewidth-2\tabcolsep-2\arrayrulewidth\relax}|}{1. BRUCE, Alberts et al. \textbf{Biologia molecular da célula}. 6. ed. Porto Alegre: Artmed, 2017.
2. CAMPBELL, Mary K.; FARRELL, Shawn O. \textbf{Bioquímica}. 2. ed. São Paulo: Cengage Learning, 2016.
3. KIERSZENBAUM, Abraham L.; TRES, Laura L. \textbf{Histologia e biologia celular: uma introdução à patologia}. 3. ed. Rio de Janeiro: Elsevier, 2012.} \\ \hline
\end{xltabular}

\vspace{0.4cm}


% ---------------------------------------------------------
% TABELA EMENTA 13: Biologia — Ano 3
% ---------------------------------------------------------
\begin{xltabular}{\linewidth}{|X|p{2.5cm}|p{2.5cm}|}
\hline
\multirow{3}{=}{\textcolor{ifscgreen}{\textbf{Unidade Curricular:}}\\[3pt]\textbf{\large Biologia — Ano 3}} & \multicolumn{2}{p{\dimexpr5cm+2\tabcolsep+\arrayrulewidth\relax}|}{\textcolor{ifscgreen}{\textbf{Semestre:}} \textbf{1º e 2º}} \\ \cline{2-3}
 & \textcolor{ifscgreen}{\textbf{CH EaD*:}} & \textcolor{ifscgreen}{\textbf{CH Total*:}} \\ \cline{2-3}
 & \textbf{00 h} & \textbf{80 h} \\ \hline
\endhead
\multicolumn{3}{|p{\dimexpr\linewidth-2\tabcolsep-2\arrayrulewidth\relax}|}{\textcolor{ifscgreen}{\textbf{Objetivos:}}} \\ \hline
\multicolumn{3}{|p{\dimexpr\linewidth-2\tabcolsep-2\arrayrulewidth\relax}|}{
\begin{itemize}[leftmargin=*,noitemsep,topsep=2pt]
 \item Analisar a estrutura e a expressão do material genético, compreendendo os mecanismos da variabilidade genética e da hereditariedade e suas aplicações em diferentes contextos científicos e tecnológicos.
 \item Avaliar criticamente os avanços da Biotecnologia e suas implicações éticas, sociais, ambientais e econômicas, fundamentando-se em evidências científicas para a tomada de decisões responsáveis frente aos desafios tecnológicos e socioambientais contemporâneos.
 \item Compreender a diversidade dos seres vivos, reconhecendo suas categorias taxonômicas, relações evolutivas, características biológicas, estratégias adaptativas, interações com o ambiente e importância para o equilíbrio dos ecossistemas e a sustentabilidade das sociedades humanas.
 \item Desenvolver habilidades de observação e experimentação científica de modo a identificar e explicar fenômenos biológicos e processos tecnológicos decorrentes da utilização de organismos vivos na produção industrial, na saúde pública e na conservação ambiental.
 \item Valorizar a diversidade biológica e cultural dos ecossistemas regionais, reconhecendo a contribuição dos saberes tradicionais africanos e indígenas para a relação com a natureza, o uso sustentável dos recursos naturais e a promoção da saúde individual e coletiva.
\end{itemize}
} \\ \hline
\multicolumn{3}{|p{\dimexpr\linewidth-2\tabcolsep-2\arrayrulewidth\relax}|}{\textcolor{ifscgreen}{\textbf{Conteúdos:}}} \\ \hline
\multicolumn{3}{|p{\dimexpr\linewidth-2\tabcolsep-2\arrayrulewidth\relax}|}{
\begin{itemize}[leftmargin=*,noitemsep,topsep=2pt]
 \item Biologia molecular: DNA, RNA e expressão gênica (transcrição e tradução).
 \item Genética e hereditariedade: conceitos básicos, leis de Mendel e padrões de herança.
 \item Temas contemporâneos em Biotecnologia e Bioética: células-tronco, clonagem, transgênicos, edição genética, experimentação animal, genética forense, entre outros.
 \item Evolução biológica: teorias, evidências e mecanismos evolutivos.
 \item Diversidade e classificação dos seres vivos: vírus, bactérias, algas, protozoários, fungos, plantas e animais.
 \item Ecologia: sistemas ecológicos e serviços ecossistêmicos.
 \item Educação Ambiental e Cultura Oceânica.
 \item História e Cultura Afro-Brasileira, Africana e Indígena.
\end{itemize}
} \\ \hline
\multicolumn{3}{|p{\dimexpr\linewidth-2\tabcolsep-2\arrayrulewidth\relax}|}{\textcolor{ifscgreen}{\textbf{Estratégias de Ensino e Aprendizagem:}}} \\ \hline
\multicolumn{3}{|p{\dimexpr\linewidth-2\tabcolsep-2\arrayrulewidth\relax}|}{- **Metodologia:** Abordagem dialógica, contextualizada e investigativa, com levantamento de saberes prévios e problematização de situações reais. Aulas expositivo-dialogadas com quadro e projetor multimídia; atividades práticas em campo e laboratório; análise de experimentos e estudos de caso; discussão de textos, imagens e vídeos científicos; uso de atlas digitais, softwares de animação e simulação, modelos didáticos tridimensionais e jogos educativos; debate de temas transversais; visitas técnicas e eventos científicos. Aulas práticas realizadas preferencialmente no pátio do câmpus e no Laboratório de Biociências (média de 20h anuais). SIGAA utilizado como Ambiente Virtual de Aprendizagem.
- **Avaliação:** Perspectiva processual e formativa, observando a evolução dos conhecimentos ao longo do período letivo, por meio de estudos dirigidos; listas de exercícios focados em ENEM e vestibulares; provas teóricas; relatórios de atividades experimentais; projetos, seminários ou oficinas temáticas; e avaliação atitudinal.} \\ \hline
\multicolumn{3}{|p{\dimexpr\linewidth-2\tabcolsep-2\arrayrulewidth\relax}|}{\textcolor{ifscgreen}{\textbf{Bibliografia Básica:}}} \\ \hline
\multicolumn{3}{|p{\dimexpr\linewidth-2\tabcolsep-2\arrayrulewidth\relax}|}{1. GRIFFITHS, Anthony J. F. \textbf{Introdução à genética}. Tradução de Idilia Vanzellotti. 10. ed. Rio de Janeiro: Guanabara Koogan, 2013.
2. REECE, Jane B. \textbf{Biologia de Campbell}. 10. ed. Porto Alegre: Artmed, 2015.
3. Livro didático fornecido pelo Fundo Nacional de Desenvolvimento da Educação (FNDE).} \\ \hline
\multicolumn{3}{|p{\dimexpr\linewidth-2\tabcolsep-2\arrayrulewidth\relax}|}{\textcolor{ifscgreen}{\textbf{Bibliografia Complementar:}}} \\ \hline
\multicolumn{3}{|p{\dimexpr\linewidth-2\tabcolsep-2\arrayrulewidth\relax}|}{1. BEGON, Michael; TOWNSEND, Colin R.; HARPER, John L. \textbf{Ecologia: de indivíduos a ecossistemas}. 4. ed. Porto Alegre: Artmed, 2007.
2. CASTRO, Peter; HUBER, Michael E. \textbf{Biologia marinha}. 8. ed. Porto Alegre: Artmed, 2012.
3. DARWIN, Charles. \textbf{A origem das espécies e a seleção natural}. Tradução de Soraya Freitas. São Paulo: Madras, 2011.} \\ \hline
\end{xltabular}

\vspace{0.4cm}

